\section{Analysis strategy and the pre-registered working-point rule (P2)}\label{sec:strategy}

The finder must select quirk planes at a pile-up fake rate of order $10^{-9}$ per triggered event, far below what can be
measured directly. The strategy therefore separates three steps that are fixed in this order: (i) a finder and an ordered
ladder of selections defined on background samples only; (ii) a rule, frozen before any signal efficiency was computed,
that picks the working point (WP) from the background estimates; (iii) a pre-defined fallback deliverable if no rung
satisfies the rule. This section defines all three; the frozen text is \rep{acts/stage4/WP_RULE.md} (sha256
\texttt{1b2ea9bd\ldots}) with the amendment \rep{WP_RULE_PC19.md}, and every later change is logged in
\rep{acts/stage4/HASH_ADDENDUM.md}.

\subsection{Trigger and offline preselection}\label{sec:strategy:trig}
The quirks are nearly invisible to the calorimeters, so the event is triggered by the missing transverse momentum (MET)
from initial-state radiation, as in previous proposals~\cite{Knapen:2017kly,Li:2019wce}; the Level-1 track trigger is
not used, because quirk hits form few stubs and no helices. The trigger is modelled as an efficiency on the smeared
MET: $\Eff_{\rm trig}=0.98\cdot\frac12\{1+\operatorname{erf}[(\mathrm{MET}-200\GeV)/(\sqrt2\cdot20\GeV)]\}$ (\EST{};
the Phase-2 L1 PUPPI-MET thresholds~\cite{Zabi:2020gjd} are carried as $\pm25\GeV$ on the turn-on point). The offline
preselection requires MET$\,>250\GeV$, at least one jet with $p_T>100\GeV$ and $|\eta|<2.4$, $\min\Delta\phi(\text{jet}_{1-4},
\text{MET})>0.5$, and vetoes isolated electrons and muons ($p_T>10\GeV$), hadronic $\tau$ ($p_T>20\GeV$) and photons
($p_T>15\GeV$); the vetoes cost the signal nothing and act only on the SM backgrounds. The Stage-3 signal samples
contain only pairs with $p_T(Q\bar Q)>200\GeV$ (and $|\eta_{\rm pair}|<4$), the recoil tail that the trigger selects, so
the yield carries the tail fraction $f_{200}$ of the inclusive cross section (Table~\ref{tab:xsec} for $E$: 0.048, 0.104, 0.146
and 0.169 at 200, 500, 1000 and $1500\GeV$; 0.114 for $\Etil$ at $500\GeV$). The signal efficiency factorises into
$f_{200}$, a truth-level acceptance $A$ of the tail sample and a detector-level selection efficiency $\Eff_{\rm sel}$:
\begin{equation}
  N_{\rm exp}=\sigma_{\rm norm}\,f_{200}\,L\,A,\quad A=\Eff_{\rm fid}\,\Eff_{\rm trig}\,\Eff_{\rm presel|trig}\,\Eff_{\rm PV},
  \quad S=N_{\rm exp}\,\Eff_{\rm sel},\quad \Eff_{\rm tot}=f_{200}\,A\,\Eff_{\rm sel},
  \label{eq:effchain}
\end{equation}
with $L=3\abinv$ and the inclusive $\sigma_{\rm norm}$ of Eq.~\eqref{eq:signorm}. $\Eff_{\rm fid}$ is the fraction of tail
pairs with $|\eta_{\rm pair}|<2.5$, $\Eff_{\rm presel|trig}$ the preselection efficiency after the trigger and
$\Eff_{\rm PV}=1$ (PV identification not modelled); $A_{\rm incl}=f_{200}A$ is the acceptance relative to the inclusive
cross section (Table~\ref{tab:acc}). $N_{\rm exp}$ is the predicted signal yield in the acceptance, $\Eff_{\rm sel}$ the
fraction of these events with a signal-matched candidate that passes the selection at pile-up 200 and has
$f_{\rm CKF}<0.2$ (events whose quirks did not both leave the tracker in the truth simulation count as failures), and $S$
the expected number of selected signal events; the factorisation assumes $\Eff_{\rm sel}$ independent of the
preselection. From Section~\ref{sec:results} on, $\Eff$ without subscript means $\Eff_{\rm sel}$. A candidate counts as signal-matched only if at least three of its core hits on at least three layers
are quirk hits and quirk hits are at least half of its core hits (declared before any signal efficiency was looked at).

\subsection{The plane finder}\label{sec:strategy:finder}
The baseline finder (F0) looks for a plane through the primary vertex (PV) that contains repeated crossings of the same
layers, in the azimuthal road of the MET. In PV coordinates $\mathbf u=\mathbf x-\mathbf x_{\rm PV}$:
\begin{enumerate}
\item \textbf{Road and window.} Hits must satisfy $|\Delta\phi(\mathbf u,\phi_{\rm MET})|<\dninetyfive+\arcsin\min(1,W_c/r)$,
  lie forward of the PV along the plane axis and within $|\mathbf u\cdot\hat n_2|<W_c$ of it, with $W_c=75\mm$ (half the
  outer IT radius, rounded up, from $\lmax\le\RIT$). The road half-width $\dninetyfive=0.15$\,rad is the 95\% quantile
  of $|\Delta\phi(\mathrm{MET},p_T(Q\bar Q))|$ at generator level after trigger and preselection, maximised over the 17
  (species, mass) points (no detector response; 0.130--0.149\,rad per mass).
\item \textbf{Seeds.} Pairs of hits on the outer two IT barrel layers with $|\Delta\phi|<0.3$, $|\Delta\eta|<0.2$, confirmed
  by a third hit on the second layer, define a plane normal $\hat n$; a seed is refitted if at least $\Tseed$ layers have
  $\ge2$ crossings in its band. $\Tseed=6$ was frozen and changed to $\Tseed=4$ by the action pre-declared for the G3 pilot
  (Section~\ref{sec:bkg:closure}).
\item \textbf{Collection and refit.} Hits with $|\hat n\cdot\mathbf u|<3\sigma_{\rm tot}$, where
  $\sigma_{\rm tot}^2=\sigma_n^2+t_{\rm tol}^2$, $\sigma_n$ is the hit resolution projected on $\hat n$ and
  $t_{\rm tol}=30\um$ is the plane-thickness tolerance (the truth-level B+MS thickness is 10--30$\um$ above
  $\Lcross\approx1.2$--$1.9\keV$, Section~\ref{sec:p1:ms}). Three iterations of a PV-constrained weighted plane fit follow; duplicate
  candidates within 2\,mrad are merged.
\item \textbf{Variables.} Per-layer crossing counts (on stacked OT modules the minimum of the two sensors): $N_{\rm IT2}$, $N_{\rm
  IP2}$ and $N_{\rm 2S2}$ count IT, IT+PS and 2S barrel layers with at least two crossings. The robust core width
  $\wcore=1.4826\,\mathrm{MAD}(\mathbf u\cdot\hat n_2)$ over the in-band IT and PS hits measures the ribbon width (a proxy of
  $\lmax$), and $\chicore$ is the plane-fit $\chi^2$ of the core hits with its $p$-value. $f_{\rm CKF}$ is the fraction of
  band hits that sit on CKF tracks.
\item \textbf{Fixed cuts at every rung.} $N_{\rm IT2}\ge3$, $|\Delta\phi(\hat n_3,\mathrm{MET})|<\dninetyfive$,
  $|\eta(\hat n_3)|<1.0$ and $f_{\rm CKF}<0.2$ (the last not credited for the pile-up background). The pseudorapidity cut on
  the plane axis $\hat n_3$ restricts the analysis to the barrel (user decision, Appendix~\ref{app:limits}); 1.0 is the edge of the third 2S
  layer (0.94) rounded up, a geometry-only number.
\end{enumerate}
The improved variant (F1) adds three quirk-hypothesis variables computed from the stored hits: a helix test $h$ (with the
PV and with a free origin), the linearity $o$ of the opening of the two arms, and a truncated-mean $dE/dx$ in MIP units, the pooled truncated mean of all
muon-gun TBPX hits ($3.558\times10^{-4}\GeV$). Candidates are split into two categories: \textbf{SR-T},
in which the earliest in-band BTL hit is late by $\Delta t=t_{\rm BTL}-t_{\rm PV}-|\mathbf u|\in[t_{\min},t_{\max}]$
(natural units), where $t_{\min}$ is set where the per-hit rate of in-time pile-up falls to $10^{-4}$ and
$t_{\max}$ where that of the next bunch crossing rises to $10^{-4}$, both from pile-up-only events (logged amendment, Appendix~\ref{app:limits}; $\sigma_{t,\rm PV}=190\ps$, varied to 150 and 230\,ps), and \textbf{SR-G}, all others.
The binding window, from the production pile-up-only sample, is $[5.955,\,24.035]\ns$, and the
pre-registered time-independence test passes.

\subsection{The ladder}\label{sec:strategy:ladder}
Each category walks its own copy of a ladder of 23 cumulative rungs (Table~\ref{tab:ladder}); each rung tightens one
threshold of the previous one. The order is fixed by background considerations only: counting rungs first (the base of
the pile-up extrapolation), then width, core $\chi^2$, 2S layers and the F1 variables. There is no fixed cut on the plane
width; the width thresholds are rungs.

\begin{table}[htbp]
\centering
\small
\caption{The frozen ladder (\rep{WP_RULE.md} \S3), applied with the fixed cuts of Section~\ref{sec:strategy:finder}.
Each rung adds the stated requirement to all previous ones. R04 ($N_{\rm IP2}\ge7$) means a double crossing on every
IT and PS barrel layer.}
\label{tab:ladder}
\begin{tabular}{lll}
\toprule
rungs & family & requirement added\\
\midrule
R01--R04 & F0 counting & $N_{\rm IT2}\ge3$ and $N_{\rm IP2}\ge4,5,6,7$\\
R05--R08 & F0 width & $\wcore<37.5,\,18.75,\,9.4,\,4.7\mm$ ($W_c/2^k$, $k=1\ldots4$)\\
R09--R11 & F0 core $\chi^2$ & $p(\chicore)>10^{-3},\,0.01,\,0.05$\\
R12--R14 & F0 2S & $N_{\rm 2S2}\ge1,2,3$ (R14 = full F0)\\
R15--R17 & F1 helix & $h>2,3,5$\\
R18--R20 & F1 opening & $o>2,3,5$\\
R21--R23 & F1 $dE/dx$ & $dE/dx>1.3,\,1.6,\,2.0$ MIP\\
\bottomrule
\end{tabular}
\end{table}

\subsection{Background budget and the working-point rule}\label{sec:strategy:rule}
For each category $\kappa\in\{\text{T},\text{G}\}$ (SR-T, SR-G) and rung $r$ the expected background in $3\abinv$ is the sum
\begin{equation}
  \Nbkg = \Ntrig\Big[P_{\rm PU}+P_{\rm SM,HS}+\sum_t \text{rate}_t\,P(\text{pass}\mid t)\Big]+N_{\rm cosmic}+N_{\rm halo},
  \label{eq:bsum}
\end{equation}
where $\Ntrig$ is the number of SM events after trigger and preselection, $P_{\rm PU}$ the pile-up-only fake probability
per event, $P_{\rm SM,HS}$ that of candidates involving hard-scatter hits, and $t$ runs over single-interaction
topologies (conversions, $V^0$, PV pairs, loopers, nuclear interactions, decays in flight) with generator-level rates
(Section~\ref{sec:bkg}). Each term is a product of measured factors with a log-normal uncertainty,
$\sigma_{\ln}^2=\text{stat}^2+(\text{extrapolation model})^2+(\text{density})^2+(\text{normalisation})^2$, and
$\Nbkgeightyfour$ is the 84th percentile of the Monte-Carlo sum. A factor with zero observed counts enters once per term,
at $1.14/M$ (the one-sided 68\% Poisson bound for a population of size $M$) with $\sigma_{\ln}=1$, never as zero.

\paragraph{Rule.} The per-category WP is the loosest rung $r^*_\kappa$ with $\Nbkgeightyfour\le0.11$ events for it and every
tighter rung. The value 0.11 follows from log-normal algebra alone: two independent categories at $\sigma_{\ln}=2.3$
combine at 2.71 times the per-category $\Nbkgeightyfour$, and $0.3/2.71=0.11$. If both categories have a WP, the pair is
accepted only if the combined $\Nbkgeightyfour\le0.3$ for it and for every tighter pair, with a deterministic continuation
otherwise. If no rung qualifies, the category has no WP: the user is told, and the pre-defined fallback applies. No
re-ordering of the ladder and no new rung are allowed after the freeze, and further production is decided on cost and
background reach, never on signal yields.

\subsection{The fallback deliverable}\label{sec:strategy:fallback}
The fallback (WP\_RULE \S8, approved by the user before production) defines one admissible selection per
hypothesis from truth-level quantities published in Stage~3, so that a selection never cuts away the hypothesis it tests:
\begin{itemize}
\item a width rung is admissible only if $\Lam\ge\Lam_R(\mQ)\,2^{(k+1)/2}$, where $\Lam_R(\mQ)$ is the smoothed $\Lam$
  at which the median $\lmax$ equals $\RIT$ (Fig.~\ref{fig:edges}); $k_{\rm adm}\in\{0,\ldots,4\}$ is the largest such $k$;
\item a $dE/dx$ rung is admissible only if its threshold is below $R_{\rm eff}/1.25$, where $R_{\rm eff}$ is the expected
  truncated-mean ionisation of the quirk at its IT velocity, diluted by pile-up hits in the core (a frozen mixture model
  using the PDG restricted-loss formula for silicon~\cite{ParticleDataGroup:2024cfk}); this happens only at the tightest
  width rung for heavy masses;
\item the fallback selection is $r_{\rm fb}={\rm FB}_{k_{\rm adm}}\wedge h>5\wedge o>5\,(\wedge\,dE/dx>e_{\rm adm})$, plus
  the time tag in SR-T, where ${\rm FB}_k\equiv N_{\rm IT2}\ge3\wedge N_{\rm IP2}\ge7\wedge p(\chicore)>0.05\wedge
  N_{\rm 2S2}\ge3\wedge\wcore<W_c/2^k$ is the non-width part of R14 with a single width cut (FB0: no width cut);
  $e_{\rm adm}$ is the largest admissible $dE/dx$ threshold.
\end{itemize}
The primary result at each $(\mQ,\Lam,\kappa)$ is the admissible WP if it exists, otherwise the fallback selection with its
non-zero log-normal background. Choosing the best selection per point is forbidden. Efficiencies and backgrounds are also
reported at five fixed reporting selections R$_A$ (end of counting) to R$_E$ (full F1), without limits.

\paragraph{Statistical treatment.} The signal strength is the yield $\mu$ in the acceptance, with
$\mu\,\Eff_{{\rm sel},\kappa}$ expected selected events in category $\kappa$ ($\mu=N_{\rm exp}$ for the prediction).
Its expected 95\% CL upper limit $\mu_{95}$ per $(\mQ,\Lam)$ uses a two-bin
(SR-T, SR-G) counting likelihood with each background as a log-normal nuisance and the CL$_s$
method~\cite{Junk:1999kv,Read:2002hq}, evaluated with hybrid toys (5000 per point; primary) and the asymptotic
formulae~\cite{Cowan:2010js} as a cross-check, in an in-house implementation. After the last analysis fix (a
saturated-likelihood offset in the asymptotic mode) the asymptotic/toy ratio of the limits is 0.32--1.07 on the final
run: 1.02--1.07 where one category carries all the signal, 0.32--0.72 where both do and $\sigma_{\ln}$ is large, a
genuine difference between the two methods. Toys are therefore primary, and the asymptotic $\pm1\sigma$ bands are not
used. A comparison with pyhf~\cite{Heinrich:2021gyp} exists only for the dry run before that fix and is not claimed for the final
code. The limit on the inclusive cross section is $\sigma_{95}=\mu_{95}/(L\,f_{200}\,A)$, so that
$\sigma_{95}/\sigma_{\rm norm}=\mu_{95}/N_{\rm exp}$ (the figures label $\sigma_{\rm norm}$ as $\sigma_{\rm th}$); a ratio
below one would be an exclusion. Because $\mu_{95}$ is a limit on the yield, the signal-normalisation uncertainties
(the $\pm5\%$ on $\sigma_{\rm norm}$ and the tail rate of Section~\ref{sec:model:syst}) do not enter it.

\subsection{Blinding, freeze and logged changes}\label{sec:strategy:freeze}
The rule, the finder code and its configuration were hash-frozen before the first signal efficiency (the only
signal-derived inputs are generator- or truth-level quantities that the plan pre-declares: $\dninetyfive$, the IT
velocities and $\Lam_R$). Every later change is logged in \rep{acts/stage4/HASH_ADDENDUM.md} and was pre-declared or
regression-tested: runner hooks and a recovery for an ACTS navigation assertion (bit-identical on the existing job kinds);
analysis-level fixes from the analysis reviews, which changed background estimates but no selection; $\Tseed=6\to4$
(G3 pilot, Section~\ref{sec:bkg:closure}); the SR-T window, because the frozen window $[t_{\min},25\ns]$ was empty once
next-crossing hits were included; and the full-chain re-run (A2FC) of nine signal points under the pre-registered revert
rule.
